\section{Data and tasks}
\label{sec:data}

Four task families cover the settings studied here. Three of them impute
within one survey. There the receivers' true values are known, so we split
the survey into donors and receivers by the protocol of
Section~\ref{sec:protocol} (80/20 holdouts repeated over ten paired seeds)
and score each method against the truth. The fourth, the population-view
experiment, imputes onto a second survey, where no true values exist, and
scores each candidate file against a held-out part of the donor survey.
Every table in Section~\ref{sec:results} regenerates from the repository's
committed run configurations; row limits and skipped cells are recorded in
the run manifests.

\subsection{Task 1: wealth within the SCF}
\label{sec:task-scf}

The Survey of Consumer Finances \citep{fed2023scf} observes household wealth
in detail. From the 2022 summary extract (first implicate, survey weights) we
impute debt (zero for many households, never negative) and net worth
(negative, zero, or positive, with a heavy right tail) from age, sex,
education class, marital status, children, total income, and wage income.
The two targets are chosen deliberately: their sign structures exercise the
sign gates, and their mutual dependence exercises chaining.
Appendix~\ref{app:scf2025} reruns this task and both profiles of the next
on the 2025 wave, which no method here was tuned on.

\subsection{Task 2: the SCF\texorpdfstring{$\to$}{->}CPS population view}
\label{sec:task-harness}

Each method fits wealth conditionals on an SCF donor split and imputes them
onto real households from the \cps{}'s Annual Social and Economic Supplement
(ASEC) that share the predictor set. The households are built from the
Census Bureau's ASEC public-use files \citep{census2025}, read directly from
the Bureau's website so that no task input flows through any PolicyEngine
processing. The experiment runs in two
profiles. The \emph{minimal} profile has six shared predictors, the two
wealth targets of Section~\ref{sec:task-scf}, and a single receiver vintage
(ASEC 2025). It is the controlled case, small enough to trace each
mechanism. The \emph{production-scale} profile mirrors the production
stack's shape in three ways. The shared predictor set widens to the ten
variables the two surveys share, adding education class, race,
homeownership, and labor-force status, each mapping verified against both
files' codebooks. The receiver pools three ASEC vintages (2023--2025,
weights scaled to one household population, 168,852 households), as
\populace{}'s production build does. And the targets widen to a four-component
chained balance sheet (financial assets, nonfinancial assets, debt, then
net worth), whose accounting identity
$\text{networth} = \text{fin} + \text{nfin} - \text{debt}$ holds exactly in
the extract, so chaining's contribution can be read against the identity. In both
profiles each run imputes onto a fixed uniform subsample of 20,000 receiver
households (of 55,762 in the minimal profile and 168,852 at production
scale), the same for every method and seed. The candidate population
(\cps{} demographics plus imputed wealth, under \cps{} household weights)
is then scored through the harness of Section~\ref{sec:harness} against the
held-out SCF, on shared predictors and wealth together, under SCF
weights. The sampling-floor row of each results table scores the donor split
itself against the holdout. This
experiment compares how the methods reproduce the joint distribution of
shared predictors and wealth on a receiver file the donor survey never saw. The
conditional independence that statistical matching assumes is one way to
lose that structure \citep{rassler2002statistical, dorazio2006statistical}.
Dependence between wealth and variables that only the \cps{} observes is not
scored, and every method relies on conditional independence there. Because
the demographic part of every candidate is the same real \cps{} sample,
differences between methods come from the imputed wealth and from how it
varies with demographics. The difference between the two surveys'
demographics is common to all methods.

\subsection{Task 3: zero-inflated income components in the CPS}
\label{sec:task-cps}

From the Census Bureau's ASEC 2025 public-use person file (income reference
year 2024), we impute interest and dividend income for adults from age, sex,
and employment income, under the ASEC person supplement weights. The file is
read directly from the Census Bureau's website and never from a processed or
enhanced file, which would embed the production imputations this paper
is about. The task uses the survey's published person weights. Every method
shares them, so the comparison between methods does not depend on that
choice.

Interest is zero for 41.7 percent of adults and dividends for 84.7 percent
(weighted shares), and both have heavy positive tails, which is the
structure the gates are built for. The task runs on a fixed uniform subsample of 20,000
of the file's 108,044 adults. Public-use ASEC amounts are topcoded, which truncates the
extreme tail relative to an administrative or SCF-style source. The fragility
values of Section~\ref{sec:results-cps} are those of a topcoded file.
This task also carries the reweighting stress test of
Section~\ref{sec:metrics}. The test is motivated by prior \policyengine{}
production experience, in which an unweighted fit broadcast rare high-value
donor records across a receiver file as near-zero-weight point masses that
later reweighting amplified. All fragility numbers reported here are
measured under this paper's protocol.

\subsection{Task 4: cross-dataset transfer}
\label{sec:task-openml}

To test behavior beyond economic microdata, we run the methods on
the six OpenML AutoML Benchmark regression datasets used in the
\microimpute{} manuscript's appendix \citep{vanschoren2014openml}:
\texttt{space\_ga}, \texttt{elevators}, \texttt{brazilian\_houses},
\texttt{onlinenewspopularity}, \texttt{abalone}, and \texttt{house\_sales}.
Each dataset's target is the imputed column, under uniform weights. The two
datasets with more than 20,000 rows, \texttt{house\_sales}
and \texttt{onlinenewspopularity}, are capped at 20,000 by a fixed uniform
subsample.